\documentclass[10pt,a4paper]{amsart}
\usepackage[margin=19mm]{geometry}
\usepackage{amsmath,amssymb,mathtools}
\usepackage[scr=rsfs,cal=euler]{mathalfa}
\usepackage{xfrac}
\usepackage[unicode,hidelinks,bookmarks=false]{hyperref}
\newcommand{\ie}{i.e.\ }
\newcommand{\cf}{cf.\ }
\newcommand{\resp}{respectively\ }
\newcommand{\Subsection}{\subsection}
\providecommand{\nobreakdash}{\nobreak\hbox{-}}
\setlength{\emergencystretch}{2em}
\multlinegap0pt
\usepackage{bm}
\usepackage{bbm}
\usepackage{dsfont}
\usepackage{xspace}
\usepackage{cancel}
\usepackage{mathtools}
\usepackage{amsthm}
\makeatletter
\@ifundefined{thm}{\newtheorem{thm}{Thm}}{}
\@ifundefined{claim}{\newtheorem{claim}[thm]{Claim}}{}
\makeatother
\newcommand{\SymGroup}[1][n]{\mathfrak{S}_{#1}}
\title[Configuration spaces on a wedge of spheres and Hochschild-Pi]{Configuration spaces on a wedge of spheres and Hochschild-Pirashvili homology}
\author{Nir Gadish, Louis Hainaut}
\date{Source: arXiv:2202.12494v2 (CC BY 4.0)}
\begin{document}
\maketitle

\noindent\emph{Source and licence.} Configuration spaces on a wedge of spheres and Hochschild-Pirashvili homology. Version arXiv:2202.12494v2 (CC BY 4.0), distributed under the Creative Commons Attribution 4.0 International licence, as recorded in the arXiv licence metadata for this exact version and shown on the abstract page https://arxiv.org/abs/2202.12494v2. Full rights notice: licenses/arxiv-2202.12494-CC-BY-4.0.txt.\par\smallskip

\noindent\textit{Editorial scope.} Counted unit: the Claim whose source label is \texttt{geometric-realization-labeled} in the article (source0.tex lines 846--853, source label \texttt{geometric-realization-labeled}), delivered together with the article's own complete original proof of it (source0.tex lines 855--898, one \texttt{proof} environment of 44 lines). No mathematics is written, repaired or re-derived: statement and proof are the article's own text line for line. Required context retained uncounted: none. Every definition, display and label of those lines is retained unchanged. Omitted from this package and not redistributed: 1--845, 854--854, 899--2695. Nothing inside a retained excerpt is omitted or altered: every retained body is delivered exactly as the article's lines are written, so no omission receipt of any kind accompanies this package. Citations: the retained excerpts carry no citation command at all, so no bibliography entry of the article is transcribed and no reference is redistributed. Reference adaptation: none was needed and none was made; every reference inside the delivered unit resolves inside it, so no unresolved reference or citation is produced. Printed slips kept literal: no printed word, symbol, hypothesis or number of the counted statement or of its proof was altered. Layout and class: the article's own class is \texttt{article} with packages including geometry, inputenc, fontenc, amsmath, amsthm, amssymb, array, stmaryrd, graphicx, biblatex, hyperref, url, bbm, bm; the wrapper is the lane's pinned \texttt{amsart} editorial wrapper at 10pt on A4, which restates the theorem-like environments the retained text needs and adds the article's own macro definitions that the retained text needs (\texttt{\textbackslash SymGroup}), with extra wrapper packages (bm, bbm, dsfont, xspace, cancel, mathtools). The delivered numbering is the unit's own. Assets: no retained span contains a figure environment, a graphics-inclusion command, a PDF-page inclusion, a table or a TikZ picture; no image file is copied into this package, the package declares no asset dependency and no image file is redistributed with it. Licence: version arXiv:2202.12494v2 of this article is distributed under the Creative Commons Attribution 4.0 International licence, as recorded in the arXiv licence metadata for this exact version and shown on the abstract page https://arxiv.org/abs/2202.12494v2. A full rights notice is delivered at licenses/arxiv-2202.12494-CC-BY-4.0.txt. Characters: every retained excerpt is pure ASCII, so the delivered engine, XeLaTeX with its default Latin Modern text font, reports no missing character.
\begin{claim}\label{geometric-realization-labeled}
Let $X_+$ be a simplicial set with a disjoint basepoint and finitely many non-degenerate simplices. Given a finite set $T$, the configuration space with labels in the square-zero monoid $M[T]$ has geometric realization 
\begin{equation*}
%|\ConfLabeled{(X_+)}{Z[S]}| 
|M[T]_{\bullet}\circ (X_+)| \cong  \bigvee_{n\geq 0} \left(F(|X|,n)\times_{\SymGroup[n]} T^n   \right)^+
\end{equation*}
where $Y^+$ is the one-point compactification of a topological space $Y$. That is, the resulting labeled configuration space is the wedge sum of compactified configuration spaces of particles in $|X|$ with decorations in $T$.
\end{claim}

\begin{proof}

First, we prove that there is an isomorphism of simplicial sets 
\[
    M[T]_{\bullet}\circ (X_+) \cong \bigvee_{n\geq 0}{\Big(X^n/{\Delta_n(X)}\Big) \wedge_{\SymGroup}(T^n)_+}
\]
where $\Delta_n(X) \subseteq X^n$ is the \emph{fat diagonal} -- the simplicial subset whose $p$-simplices are the $n$-tuples of elements in $X_p$ not all of whose coordinates are distinct; equivalently its $p$-simplices are the non-injective functions from $\bm{n}$ to $X_p$.% Simplicial maps of $X$ induce simplicial maps on $\Delta_n(X)$ by their action on each coordinate.

On the left-hand side, $p$-simplices are described as follows. There is a natural bijection $M[T]_{\bullet}((X_p)_+) \cong \left((\{1\}\coprod T)^{X_p}\right)_+$, further simplified by
\begin{equation}\label{eq:from functor to configurations}
\big(\{1\}\coprod T\big)^{X_p} \cong \coprod_{Y\subseteq X_p} T^Y \cong \coprod_{n\geq 0}\operatorname{Inj}(\bm{n},X_p) \times_{\SymGroup} T^n,
\end{equation}
where $\operatorname{Inj}(\bm{n},X_p)$ denotes the set of injective functions from $\bm{n}$ to $X_p$. Remembering the basepoint, we obtain the natural bijection
\[
    M[T]_{\bullet}((X_p)_+) \cong \bigvee_{n\geq 0}\operatorname{Inj}(\bm{n},X_p)_+\wedge_{\SymGroup}(T^n)_+ \cong \bigvee_{n\geq 0}\Big(X_p^n/{\Delta_n(X_p)}\Big)\wedge_{\SymGroup}(T^n)_+,
\]
where the last isomorphism uses the obvious identification $X_p^n/\Delta_n(X_p) \cong \operatorname{Inj}(\bm{n},X_p)_+$
which sends noninjective functions in $X_p^n$ to the basepoint. It remains to verify that the simplicial maps on both sides commute with these bijections. Crucially, this uses the fact that the multiplication on $T$ is trivial.

%is the basepoint if and only if some element $y\in X_q$ is labeled by $0$, which occurs if and only if in the fiber $\alpha^{-1}(y)$ more than one element is labeled by some element of $S$. In that case the image of $\phi$ under the bijection is a labeled tuple containing two elements in the same fiber of $\alpha$, and therefore on that side $\alpha_*$ also sends this element to the basepoint. Otherwise if $\alpha_*(\varphi)$ is not the basepoint, then each fiber of $\alpha$ contains no element labeled by $0$ (since $\phi$ is not the basepoint) and at most one element labeled by an element of $S$; in that case it is easy to verify that the simplicial maps $\alpha_*$ on both sides commute with the bijections. \textcolor{red}{My explanation seems far too long. You are very welcome to cut through it if you see how to make it more concise.}

Face and degeneracy maps $\alpha: X_p\to X_q$ act on $(\{1\}\coprod T)^{X_p}$ by multiplying the labels of points in every fiber (with the empty product being $1$). But since the product of two elements from $T$ is trivial, a function $\varphi\in (\{1\}\coprod T)^{X_p}$ is sent to the basepoint unless every fiber of $\alpha$ contains at most one element with label in $T$ and every other element is labeled by the unit $1$, in which case it is sent to the composition 
\[
\alpha(\varphi^{-1}(T)) \overset{\alpha^{-1}}{\longrightarrow} \varphi^{-1}(T) \overset{\varphi}\longrightarrow T
\]
extended by $1$ when this is undefined. On the RHS of \eqref{eq:from functor to configurations} the function $\varphi: \varphi^{-1}(T)\to T$ represents the element $((x_1,\varphi(x_1)),\ldots, (x_n,\varphi(x_n)))$ for some enumeration of $\varphi^{-1}(T)$. This tuple is sent under $\alpha$ to $((\alpha(x_1),\varphi(x_1)),\ldots, (\alpha(x_n),\varphi(x_n)))$, which is easily seen to be the tuple corresponding to $\alpha_*(\varphi)$ described above. 

%whenever one obtains a non-injective map. But since multiplication in $S$ is trivial, face maps either preserve injections or otherwise send them to the basepoint. Consequently, the above bijections are compatible with the simplicial maps and thus assemble into an isomorphism of simplicial sets.

%It is therefore clear that $Z[S]_\bullet\circ (X_+)$ splits as a wedge sum of the pointed simplicial sets $\left(\operatorname{Inj}(\bm{n},X)\times_{\SymGroup}S^n\right)_+$, where every non-injective function is identified with the basepoint. Note that this simplicial set is isomorphic to the one obtained by taking all functions $\operatorname{Fun}(\bm{n},X) \cong X^n$ and collapsing the simplicial subset of all non-injective functions.

Since geometric realization commutes with wedge sums, Cartesian products and quotients, one obtains a homeomorphism
\[
    |M[T]_{\bullet}\circ(X_+)| \cong \bigvee_{n\geq 0}{\Big(|X|^n/\Delta_n(|X|)\Big)\wedge_{\SymGroup}(T^n)_+},
\]
where one observes that $\Delta_n(|X|) = |\Delta_n(X)|\subseteq |X|^n$ is the topological \emph{fat diagonal}, in which at least two coordinates are equal.

Up to this point, the analysis applies to any simplicial set. When $|X|$ is compact (equivalently, $X$ has finitely many non-degenerate simplices), $|X|^n/\Delta_n(|X|)$ is homeomorphic to the one-point compactification $F(|X|, n)^+$, inducing the claimed homeomorphism.
%the same description applies after realization. Furthermore, since under the identification of geometric realizations $|X|^n \cong |X^n|$ the simplicial subset of non-injective functions is realized as the `fat diagonal' $\Delta^n(|X|)$, there is a homeomorphism of pointed spaces
%\[
%|\operatorname{Inj}(\bm{n},X)_+| \cong |X|^n/\Delta^n(|X|) \cong F(|X|,n)^+
%\]
%where the last identification follows from the equality $F(|X|,n) = |X|^n \setminus \Delta^n(|X|)$ and $|X|$ being compact. The same quotient description applies after multiplying the two spaces by $S^n$ and quotienting by the diagonal $\SymGroup$-action, therefore giving the desired form.
\end{proof}

\end{document}
